\section{FHE and leveled FHE}\label{fhe-intro}

Alice has a secret $x$, and Bob has a function $f$. They want to
compute $f(x)$. Actually, Alice wants Bob to compute $f(x)$ -- but
she does not want to tell him $x$. What Alice wants is a \emph{fully
homomorphic encryption (FHE)} protocol, meaning:

\begin{enumerate}
\item
  Alice encrypts $x$ and sends Bob $\Enc(x)$.
\item
  Bob then ``applies $f$ to the ciphertext'' and obtains
  $\Enc\left( f(x) \right)$, sending it to Alice.
\item
  Alice decrypts $\Enc\left( f(x) \right)$ to learn
  $f(x)$.
\end{enumerate}

\emph{Leveled FHE} is a weaker version of FHE. Like FHE, leveled FHE
lets you perform operations on encrypted data. But unlike FHE, there
will be a limit on the number of operations you can perform before the
data must be decrypted.

Why is there a limit? Loosely speaking, the encryption procedure will
involve some sort of ``noise'' or ``error.'' As long as the error is not
too big, the message can be decoded without trouble. But each operation
on the encrypted data will cause the error to grow --- and if it grows
beyond some maximum error tolerance, the message will be lost. So there
is a limit on how many operations you can do before the error gets too
big.

As a sort of silly example, imagine your message is a whole number
between 0 and 10 (so it is one of 0, 1, 2, 3, 4, 5, 6, 7, 8, 9, 10), and
your ``encryption scheme'' encrypts the message as a real number that is
very close to the message, and decrypts a real number by ``round to the
nearest integer.'' So, the message 2 might be encrypted as 1.999832.

(You might be thinking: This is some pretty terrible cryptography,
because the message is not secure. Anyone can figure out how to round a
number, no secret key required. Yep, you are right. The full scheme
(Section \ref{lwe-crypto}) is more complicated. But it still has this
``rounding-off-errors'' feature, and that is what we want to focus on
right now.)

Now imagine that the main operation you want to perform is addition
(optionally, say, modulo 11). Well, every time you add two encrypted
numbers ($1.999832 + 2.999701 = 4.999533$), the errors add as well.
After too many operations, the error will exceed $0.5$, and the
rounding procedure will not give the right answer anymore. But as long
as you are careful not to go over the error limit, you can add
ciphertexts with confidence.

For our leveled FHE protocol, our message will be a bit (either 0 or 1)
and our operations will be the logic gates AND and NOT. Because any
logic circuit can be built out of AND and NOT gates, we will be able to
perform arbitrary calculations within the FHE encryption.

Our protocol uses a cryptosystem built from a problem called ``learning
with errors.'' ``Learning with errors'' is kind of a strange name; it
would make more sense to call it ``approximate linear algebra modulo
$q$.'' Anyway, we will start with the learning-with-errors problem
(Section \ref{lwe}) and how to build cryptography on top of it
(Section \ref{lwe-crypto}) before we get back to leveled FHE.
